\begin{table}[H]
\centering
\captionsetup{justification=centering}
\caption{Correlação estimada entre os distúrbios das inclinações - modelo multivariado - total de desocupados}
\label{tab:matriz_correlacao_desoc}

\renewcommand{\arraystretch}{0.8}

\resizebox{\linewidth}{!}{%
\begin{tabular}{@{}p{2.8cm}c*{8}{c}@{}}
\toprule
\textbf{Estrato Geográfico} & \(\hat{\sigma}_R^2\) & \textbf{1} & \textbf{2} & \textbf{3} & \textbf{4} & \textbf{5} & \textbf{6} & \textbf{7} & \textbf{8} \\
\midrule
01 - Belo Horizonte & 33,8915 & 1 &  &  &  &  &  &  &  \\
\raisebox{0em}{\parbox[t]{2.8cm}{02 - Entorno e Colar \\[-2pt] Metropolitano de BH}} & 112,5743 & 0,9911 & 1 &  &  &  &  &  &  \\
03 - Sul de Minas & 19,8207 & 0,9801 & 0,9910 & 1 &  &  &  &  &  \\
04 - Triângulo Mineiro & 13,5033 & 0,9879 & 0,9977 & 0,9978 & 1 &  &  &  &  \\
05 - Zona da Mata & 10,7285 & 0,5403 & 0,6164 & 0,6914 & 0,6554 & 1 &  &  &  \\
06 - Norte de Minas & 27,4822 & 0,9800 & 0,9930 & 0,9997 & 0,9986 & 0,6916 & 1 &  &  \\
07 - Vale do Rio Doce & 30,7076 & 0,9527 & 0,9797 & 0,9923 & 0,9882 & 0,7555 & 0,9937 & 1 &  \\
08 - Central & 15,3830 & 0,9894 & 0,9762 & 0,9503 & 0,9654 & 0,4365 & 0,9513 & 0,9181 & 1 \\
\bottomrule
\end{tabular}%
}

\fonte{Elaboração própria, com base nos dados da PNAD Contínua. Nota: autovalores estimados de \(\Sigma_R\): 254,769; 8,858; 0,343; 0,120 e 4 valores inferiores a \(10^{-3}\). Posto numérico de 7 em 8; posto \textbf{efetivo} de 2, contando os autovalores que retêm ao menos 1\% do maior. As correlações devem, portanto, ser lidas em conjunto, e não isoladamente.}
\end{table}
